\documentclass[10pt,a4paper]{amsart}
\usepackage[margin=19mm]{geometry}
\usepackage{amsmath,amssymb,mathtools}
\usepackage[scr=rsfs,cal=euler]{mathalfa}
\usepackage{xfrac}
\usepackage[unicode,hidelinks,bookmarks=false]{hyperref}
\newcommand{\ie}{i.e.\ }
\newcommand{\cf}{cf.\ }
\newcommand{\resp}{respectively\ }
\newcommand{\Subsection}{\subsection}
\providecommand{\nobreakdash}{\nobreak\hbox{-}}
\setlength{\emergencystretch}{2em}
\multlinegap0pt
\usepackage{bm}
\usepackage{bbm}
\usepackage{dsfont}
\usepackage{xspace}
\usepackage{cancel}
\usepackage{mathtools}
\usepackage{amsthm}
\makeatletter
\@ifundefined{theorem}{\newtheorem{theorem}{Theorem}}{}
\makeatother
\newcommand{\abs}[1]{\left| #1 \right|} 
\newcommand{\derivative}[1]{#1^{\prime}}
\newcommand{\pr}[1]{ \left( #1 \right) }
\title[Numerical Analysis of HiPPO-LegS ODE for Deep State Space Mo]{Numerical Analysis of HiPPO-LegS ODE for Deep State Space Models}
\author{Jaesung R. Park, Jaewook J. Suh, Youngjoon Hong, Ernest K. Ryu}
\date{Source: arXiv:2412.08595v2 (CC BY 4.0)}
\begin{document}
\maketitle

\noindent\emph{Source and licence.} Numerical Analysis of HiPPO-LegS ODE for Deep State Space Models. Version arXiv:2412.08595v2 (CC BY 4.0), distributed under the Creative Commons Attribution 4.0 International licence, as recorded in the arXiv licence metadata for this exact version and shown on the abstract page https://arxiv.org/abs/2412.08595v2. Full rights notice: licenses/arxiv-2412.08595-CC-BY-4.0.txt.\par\smallskip

\noindent\textit{Editorial scope.} Counted unit: the Theorem whose source label is \texttt{thm:exist and unique} in the article (source0.tex lines 546--559, source label \texttt{thm:exist and unique}), delivered together with the article's own complete original proof of it (source0.tex lines 561--786, one \texttt{proof} environment of 226 lines). No mathematics is written, repaired or re-derived: statement and proof are the article's own text line for line. Required context retained uncounted: eq:legs (lines 290--293). Every definition, display and label of those lines is retained unchanged. Omitted from this package and not redistributed: 1--289, 294--545, 560--560, 787--2155. Nothing inside a retained excerpt is omitted or altered: every retained body is delivered exactly as the article's lines are written, so no omission receipt of any kind accompanies this package. Citations: the retained excerpts carry no citation command at all, so no bibliography entry of the article is transcribed and no reference is redistributed. Reference adaptation: none was needed and none was made; every reference inside the delivered unit resolves inside it, so no unresolved reference or citation is produced. Printed slips kept literal: no printed word, symbol, hypothesis or number of the counted statement or of its proof was altered. Layout and class: the article's own class is \texttt{article} with packages including arxiv, inputenc, fontenc, hyperref, url, booktabs, amsfonts, nicefrac, microtype, lipsum, graphicx, natbib, doi, stmaryrd; the wrapper is the lane's pinned \texttt{amsart} editorial wrapper at 10pt on A4, which restates the theorem-like environments the retained text needs and adds the article's own macro definitions that the retained text needs (\texttt{\textbackslash abs}, \texttt{\textbackslash derivative}, \texttt{\textbackslash pr}), with extra wrapper packages (bm, bbm, dsfont, xspace, cancel, mathtools). The delivered numbering is the unit's own. Assets: no retained span contains a figure environment, a graphics-inclusion command, a PDF-page inclusion, a table or a TikZ picture; no image file is copied into this package, the package declares no asset dependency and no image file is redistributed with it. Licence: version arXiv:2412.08595v2 of this article is distributed under the Creative Commons Attribution 4.0 International licence, as recorded in the arXiv licence metadata for this exact version and shown on the abstract page https://arxiv.org/abs/2412.08595v2. A full rights notice is delivered at licenses/arxiv-2412.08595-CC-BY-4.0.txt. Characters: every retained excerpt is pure ASCII, so the delivered engine, XeLaTeX with its default Latin Modern text font, reports no missing character.
\begin{align}   \label{eq:legs} %\tag{LegS ODE}
    \derivative{c}(t) &= -\frac{1}{t}Ac(t) + \frac{1}{t}Bf(t) 
    % c_{k+1} &= (1-\frac{A}{k})c_k + \frac{1}{k}Bf_k\\
\end{align}

\begin{theorem}[Existence and uniqueness] \label{thm:exist and unique}
% {\color{red} Shouldn't $c$ be absolutely continuous? The given solution is not necessarily differentiable.}
    For $T>0$ and $c_0 \in \mathbb{R}^N$, we say $c\colon [0,T]\rightarrow \mathbb R^N$ is a solution (in the extended sense) of the LegS ODE if $c$ is
    continuous on $[0,T]$,
    absolutely continuous on $(0,T]$, $c$ satisfies \eqref{eq:legs} for almost all $t \in (0,T]$, and $c(0) = c_0$.
    % For $M>0$, we say $c$ is a solution if $cC^0([0,M) is differentiable ;\mathbb R)$, satisfies \eqref{eq:1d_LegS} for $t \in (0,M)$ and $c(0) = c_0$.
    Assume $f\colon [0,T]\rightarrow \mathbb R$ is Lebesgue measurable, integrable, and 
    % Further, assume $f$ is continuous at $t=0$.
    % Assume locally integrable $f(t)$ 
    has a Lebesgue point at $t=0$. 
    % is continuous at $t=0$.
    Then, the solution exists and is unique if $c_0 = f(0)e_1$, where $e_1\in \mathbb{R}^N$ is the first standard basis vector.
    Otherwise, if $c_0 \ne  f(0)e_1$, a solution does not exist.
\end{theorem}

\begin{proof}
    Since $A$ is diagonalizable, the problem can be effectively reduced to $N$ $1$-dimensional problems. 
    % Using the usual spectral decomposition, let $A = VDV^{-1}$ with $D$ being a diagonal matrix. 
    Recall $A = VDV^{-1}$ where $D=\mathrm{diag}\,(1,2,\dots,N)$. 
    We see the LegS ODE \eqref{eq:legs} could be rewritten as 
    \begin{align*}
        \derivative{c}(t) 
        &= -\frac{1}{t} A c(t) + \frac{1}{t} B f(t) %\\&
        = -\frac{1}{t} VDV^{-1} c(t) + \frac{1}{t} B f(t). 
    \end{align*}
    Multiply both sides by $V^{-1}$ and denote $\tilde{c}(t) = V^{-1}c(t)$. Then, the ODE becomes
    \begin{align*}
        \derivative{\tilde{c}}(t) = V^{-1}\derivative{c}(t) &= -\frac{1}{t} D V^{-1} c(t) + \frac{1}{t} V^{-1}B f(t) = -\frac{1}{t} D\tilde{c}(t) + \frac{1}{t} V^{-1}B f(t),
    \end{align*}
    % and hence
    % \begin{align*}
    %     \derivative{\tilde{c}}(t)&= -\frac{1}{t} D\tilde{c}(t) + \frac{1}{t} V^{-1}B f(t)
    % \end{align*}
    which is a decoupled ODE with respect to $\tilde{c}$. Recalling $D_{jj} = j$, we see that the $j$-th component of the above equation is
    % Considering the $j$-th component of the above equation yields
    \begin{align}   
        \derivative{\tilde{c}}_j(t)&= -\frac{j}{t}\tilde{c}_j(t) + \frac{d_j}{t} f(t) \label{eq:1d_legs}
    \end{align}
    where $d_j = (V^{-1}B)_j$ and $\tilde{c}_j$ is $j$-th component function of $\tilde{c}$. % and noting that $(D)_{jj} = j$. 
    Since $V$ and $V^{-1}$ are absolutely continuous bijective maps from $\mathbb{R}^N$ to $\mathbb{R}^N$, 
    % the existence and uniqueness of $c(t) \in \mathbb{R}^N$ is satisfied if and only if $\tilde{c}_j(t) \in \mathbb{R}$ exists and is unique for all %$j = 1, 2, ..., N$. 
    the existence and uniqueness of the solution of the LegS ODE is satisfied if and only the existence and uniqueness of the solution of the ODE \eqref{eq:1d_legs} is satisfied for all 
    $j \in \{1, 2, ..., N\}$. 

    We now proceed by examining the existence and uniqueness of the solution 
    % $\tilde{c}_j \in C([0,T];\mathbb{R})$ 
    of the ODE \eqref{eq:1d_legs}. 
    % that is absolutely continuous on $(0,T]$, satisfying \eqref{eq:1d_legs} for almost all $t \in (0, T]$. 
    We first establish existence by presenting the explicit form of the solution.  
    Define $\tilde{c}_j\colon[0,T]\to\mathbb{R}$ as
    \begin{equation}    \label{eq:closed_form_solution}
        \tilde{c}_j(t) = \begin{cases}
            \frac{d_j}{t^j} \int^t_0 s^{j-1} f(s) ds  & \text{if} \quad t \in (0,T] \\
            \frac{d_j}{j} f(0)  & \text{if} \quad t = 0.
        \end{cases}
    \end{equation}
    By fundamental theorem of calculus, $\frac{d}{dt} \pr{ \int^t_0 s^{j-1} f(s) ds } = t^{j-1} f(t)$ holds for almost all $t\in(0,T]$ and thus $\tilde{c}_j$ is differentiable for almost all $t\in(0,T]$. 
    Therefore,
    \[
        t^j \pr{ \derivative{\tilde{c}}_j(t) + \frac{j}{t} \tilde{c}_j(t) } 
        % = t^j \derivative{\tilde{c}}_j(t) + j t^{j-1} \tilde{c}_j(t) 
        = \frac{d}{dt} (t^j\tilde{c}_j(t)) 
        = \frac{d}{dt} \pr{ d_j \int^t_0 s^{j-1} f(s) ds }
        = d_j t^{j-1} f(t) 
    \]
    holds for almost all $t\in(0,T]$. 
    Dividing both sides by $t^j$, we conclude $\tilde{c}_j$ satisfies \eqref{eq:1d_legs} for almost all $t\in(0,T]$.
    
    We now show $\tilde{c}_j$ is continuous on $[0,T]$. 
    It is sufficient to check $\tilde{c}_j$ is continuous at $t=0$ by showing
    % Since $f$ is continuous at $t=0$, from L'H\^{o}pital's rule we have
    % \begin{equation*}
    %     \begin{aligned}
    %         \lim_{t\to0^+} \tilde{c}_j(t)
    %         &= \lim_{t\to0^+} \frac{d_j}{t^j} \int^t_0 s^{j-1} f(s)ds  \\ 
    %         &= \lim_{t\to0^+} \frac{d_j }{j t^{j-1}} \pr{ t^{j-1} f(t) }
    %         = \frac{d_j}{j} \lim_{t\to0^+} f(t) 
    %         = \frac{d_j}{j} f(0) = \tilde{c}_j(0).
    %     \end{aligned}
    % \end{equation*}
    % The last equation comes from the continuity of $f$ at $t=0$. 
    \begin{align*}
        \lim_{t \to 0} \frac{d_j}{t^j} \int^t_0 s^{j-1} f(s) ds &= \frac{d_j}{j}f(0).
    \end{align*}
    % So basically we want to show that assuming $\tilde{c}_j(t)$ is 
    % \begin{equation*}
    %         \tilde{c}_j(t) = 
    %             \frac{d_j}{t^j} \int^t_0 s^{j-1} f(s) ds  \quad \text{if} \quad t \in (0,T]
    % \end{equation*}
    % that $\tilde{c}_j(0)= \frac{d_j}{j}f(0)$. Drop $d_j$ for now. 
    Since $f$ is locally integrable and has a Lebesgue point at $t=0$, observe that
    \begin{equation*}
    % \label{eq:fact_from_Lebesgue_point}
        0 = \lim_{t \to 0} 
        \frac{1}{t} \int_0^t |f(s)-f(0)|ds  
        = 
        \lim_{t \to 0} 
        \int_0^1 |f(tx) - f(0)|dx,
    \end{equation*}
    where the second equality follows by change of variables $x = s/t$. Therefore, we could deduce
    % For $t>0$, observe
    % \begin{equation*}
    %     \abs{ \tilde{c}_j(t) - \frac{1}{j} f(0) }
    %     = \abs{ \frac{1}{t} \int^t_0 \frac{1}{t^{j-1}} s^{j-1} f(s) ds - \frac{1}{t} \int_{0}^{t} \frac{1}{j} f(0) ds }
    %     \le \frac{1}{t} \int_{0}^{t} \abs{ \frac{1}{t^{j-1}} s^{j-1} f(s) - \frac{1}{j} f(0) } ds.
    % \end{equation*}
    \begin{align*}
        \abs{\lim_{t \to 0} \frac{d_j}{t^j} \int^t_0 s^{j-1} f(s) ds - \frac{d_j}{j}f(0)    
        } &\leq
        \lim_{t \to 0} 
        \frac{d_j}{t} \int_0^t \abs{ \frac{1}{t^{j-1}} s^{j-1} f(s) - \frac{1}{j} f(0)} ds \\
        &= \lim_{t \to 0} 
        d_j \int_0^1 \abs{x^{j-1} f(tx) - \frac{1}{j}f(0) } dx \\
        &\leq \lim_{t \to 0} 
        d_j \int_0^1 x^{j-1} |f(tx) - f(0)| dx  + d_j \int_0^1 \abs{x^{j-1} - \frac{1}{j}} |f(0)| dx \\
        &\leq \lim_{t \to 0} 
        d_j \int_0^1 |f(tx) - f(0)| dx  
        + d_j |f(0)| \int_0^1 \abs{x^{j-1} - \frac{1}{j} }  dx \\
        &= 0,
    \end{align*}
    % where the third inequality comes from the fact $x\le1$ and the followed equality comes from \eqref{eq:fact_from_Lebesgue_point}. 
    concluding that $\tilde{c}_j$ is continuous at $t=0$. 
    Lastly, $\tilde{c}_j$ is absolutely continuous on $(0,T]$, since for every $[t_0,t] \subset (0,T]$, 
    % $\derivative{\tilde{c}}_j$ exists almost everywhere and $\tilde{c}_j(t) = \tilde{c}_j(t_0) + \int_{t_0}^{t} \derivative{\tilde{c}}_j(s) ds$ holds.  
    both $\frac{1}{t^j}$ and $\int^t_0 s^{j-1} f(s) ds$ are absolutely continuous on $[t_0,t]$ and therefore their product is also absolutely continuous on $[t_0,t]$. 
    % \textcolor{red}{XXDRAFTSTART absolute cont at $t=0$}
    % $f$ is not absolutely continuous even if it is continuous at $t=0$. Counterexample exists. 
    % \textcolor{red}{XXDRAFTEND absolute cont}
    Hence we conclude that $\tilde{c}_j$ is a solution of the ODE \eqref{eq:1d_legs}. 

    We now establish uniqueness. 
    Suppose $\hat{c}_j$ is another solution of the ODE \eqref{eq:1d_legs}. 
    Multiplying both sides of \eqref{eq:1d_legs} by %$\exp(j\log{t}) = t^j$ 
    $t^j$ and reorganizing, for almost all $t\in(0,T]$ we have
    % \begin{align*}
    %     \frac{d}{dt} (t^jc(t)) 
    %     &= d_j t^{j-1} f(t) \\
    %     t^j c(t) - t_0^j c(t_0) &= d_j \int_{t_0}^t s^{j-1} f(s) ds
    % \end{align*}
    \[
        \frac{d}{dt} (t^j\hat{c}_j(t)) 
        = t^j \pr{ \derivative{\hat{c}}_j(t) + \frac{j}{t} \hat{c}_j(t) } 
        = d_j t^{j-1} f(t) .
    \]
    Since $\hat{c}_j$ is a solution, it is absolutely continuous on $(0,T]$, therefore $t^j\hat{c}_j(t)$ is absolutely continuous on $(0,T]$. 
    Thus for $[t_0,t] \subset (0,T]$, 
    % integrating from $t_0$ to $t$ we obtain
    by fundamental theorem of calculus we obtain
    \[
        t^j \hat{c}_j(t) - t_0^j \hat{c}_j(t_0) = d_j \int_{t_0}^t s^{j-1} f(s) ds.
    \]
    Taking limit $t_0 \to 0^+$, since $\hat{c}_j$ is a solution it is continuous at $0$, we have
    \[
        t^j \hat{c}_j(t) = d_j \int_{0}^t s^{j-1} f(s) ds.
    \]
    Dividing both sides by $t^j$ we conclude
    \begin{equation*}
       \hat{c}_j(t) = \frac{d_j}{t^j} \int^t_0 s^{j-1} f(s)ds = \tilde{c}_j(t) %\label{eq:test func}
    \end{equation*}
    for all $t \in (0,T]$. 
    % Furthermore, from L'H\^{o}pital's rule we have
    It remains to check $\hat{c}_j(0)=\tilde{c}_j(0)$. 
    Since $\hat{c}$ is continuous at $0$, we know $\hat{c}_j(0)=\lim_{t\to0^+}\hat{c}_j(t)$. Thus
     \begin{equation*}
        \hat{c}_j(0)
        = \lim_{t\to0^+} \hat{c}_j(t)
        = \lim_{t\to0^+} \tilde{c}_j(t)
        = \tilde{c}_j(0). 
    \end{equation*}
    % Again from L'H\^{o}pital's rule, we have
    % \begin{equation*}
    %     \lim_{t\to0^+} \hat{c}_j(t)
    %     % = \lim_{t\to0^+} \frac{d_j}{t^j} \int^t_0 s^{j-1} f(s)ds 
    %     % = \lim_{t\to0^+} \frac{d_j }{j t^{j-1}} \pr{ t^{j-1} f(t) }
    %     = \frac{d_j}{j} f(0) 
    % \end{equation*}
    % The last equation comes from the continuity of $f$ at $t=0$. 
    Therefore, we conclude $\hat{c}_j(t) = \tilde{c}_j(t)$ for all $t\in[0,T]$, 
     the solution of the ODE \eqref{eq:1d_legs} is unique. 
    
    % Now we continuously extend this function to $t=0$. Since $f$ is continuous at $t=0$, we can immediately apply L'H\^{o}pital's rule. 

    % Otherwise, 
    % \textcolor{blue}{Try not to overuse arrows (for limits)}
    % \begin{align*}
    %     \left| \hat{c}(t) - \frac{d_j}{j} f(0) \right|
    %     &= \left|\frac{d_j}{t^j} \int_0^t s^{j-1} f(s)ds - \frac{d_j}{j} f(0) \right| \\
    %     &= d_j \left| \frac{1}{t^j} \int_0^{t^j} \frac{1}{j}f(u^{1/j})du - \frac{1}{j}f(0) \right|
    %     % & \leq \frac{d_j}{j} \frac{1}{t^j} \int_0^{t^j} |f(u^{1/j}) - f(0)| du \\
    %     \xrightarrow[t\to 0]{}0
    % \end{align*}
    % since $t=0$ is a Lebesgue point of $f$ (\textcolor{red}{Insufficient argument}).
    
    % Hence we let $\hat{c}(0)_j = \frac{d_j}{j}f(0)$, and $\hat{c}(t)_j$ is a solution by construction. 
    % Observe that $\lim_{t\to0}\tilde{c}(t) = d_j f(0)$ since $\left| \frac{d_j}{t}\int_0^tf(s)ds - f(0) \right| \le \frac{d_j}{t}\int_0^t \left| f(s)- f(0) \right| ds \xrightarrow[t\to0]{}0$. Then by construction, this $\tilde{c}(t)$ is a solution.

    % To organize, 
    % \begin{equation*}
    %     \derivative{\tilde{c}}(t) 
    %     =  -\frac{1}{t} D\tilde{c}(t) + \frac{1}{t} V^{-1}B f(t)
    % \end{equation*}
    
    % We proceed by showing uniqueness. Assume $c_1, c_2 \in C^0[0,M)$ differentiable on $(0,M)$ satisfy (\ref{eq:1d_legs}) for $t \in (0, M)$. Then, 
    % \begin{align*}
    %     \frac{d}{dt} (c_1 - c_2)(t) &= -\frac{j}{t}(c_1-c_2)(t)\\
    %     c_1(t)-c_2(t) &= \frac{C}{t^{j}}
    % \end{align*}
    % where $C$ is a constant. However, continuity at $t=0$ enforces $C = 0$. Hence the only feasible solution is $\hat{c}(t)_j$ the one we have calculated. Moreover, this unique function has a fixed candidate for its value at $t=0$.

    As a result, we conclude the solution of the LegS ODE uniquely exists if $\tilde{c}_j(0) = \frac{d_j}{j}f(0)$, and it is given by $c = V \tilde{c}$. 
    Finally, we show the unique solution $c = V \tilde{c}$ should satisfy $c(0) = f(0)e_1$.  
    % Folding $\tilde{c}_j(0)$ back to the vector form, 
    From %from $V^{-1}c(0)_j = \tilde{c}_j(t) = \frac{d_j}{j}f(0) = \frac{1}{j}(V^{-1}B)_j f(0)$,
    \[
        V^{-1}c_j(0) = \tilde{c}_j(0) = \frac{d_j}{j}f(0) = \frac{1}{j}(V^{-1}B)_j f(0),
    \]
    % so if $c_0 =  d_j f(0)$, the solution is unique, and otherwise there is no solution.
    we see
    \[
        V^{-1}c(0)  = \text{diag} \pr{ 1, \frac{1}{2}, \frac{1}{3}, ... , \frac{1}{N} } (V^{-1}B) f(0) = D^{-1} V^{-1} B f(0).
    \]
    Multiplying both sides by $V$, we conclude
    \begin{align*}
        c(0) 
        = V D^{-1} V^{-1} B f(0)
        = ( V^{-1} D V )^{-1}Bf(0)
        % {\color{red} ??? XXX}
        = A^{-1}B f(0) 
        = f(0)e_1
    \end{align*}
    % \begin{align*}
    %     V^{-1}c(0) 
    %     &=  \text{diag}(1, 1/2, 1/3, ... ) (V^{-1}B) f(0) \\
    %     c(0) &= \underbrace{P^{-1} \text{diag}(1, 1/2, 1/3, ... ) P}_{= A^{-1}} B f(0)\\
    %     &= A^{-1}B f(0) \\
    %     &= f(0)e_1
    % \end{align*}
    % Note that this is obtained through an efficient calculation ($\mathcal{O} (N^2)$), since $A$ is a lower-triangular matrix. 
    where $e_1 = [1, 0, \dots, 0] ^{t}$.  
    Therefore if $c_0 = f(0)e_1$, the solution exists and is unique, and otherwise, there is no solution.
\end{proof}

\end{document}
